\chapter{Introdução}\label{cap:introducao}

Em novembro de 2025, Belém sediou a 30ª Conferência das Partes da ONU sobre Mudança do Clima (COP30).
O evento levou à cidade um fluxo de visitantes muito acima do normal, concentrado em poucas semanas, e
foi acompanhado de investimentos públicos em infraestrutura. Como outros megaeventos, a COP30 foi
justificada, entre outros argumentos, pela geração de empregos.

A literatura econômica trata essa promessa com cautela. Estudos \textit{ex ante}, baseados em
multiplicadores, costumam prever grandes ganhos; estudos \textit{ex post}, com dados observados,
encontram efeitos pequenos, setoriais e temporários, concentrados em hospedagem e alimentação no mês do
evento (Baade e Matheson, 2004; Hagn e Maennig, 2009; Feddersen e Maennig, 2013). Quase toda essa
evidência, porém, vem de eventos esportivos em países desenvolvidos. O desenho empírico segue os estudos que
mediram esses efeitos com registros administrativos mensais e grupos de controle formados por cidades
que não sediaram o evento (Hagn e Maennig, 2009; Feddersen e Maennig, 2013; Leoni e Boto-García, 2023).

\section{Pergunta e hipóteses}

\textbf{Pergunta de pesquisa:} o fim da COP30 esteve associado a uma redução do saldo mensal de empregos
formais em Alojamento e Alimentação em Belém e no Pará? Em que magnitude?

As hipóteses são que o saldo aumentou antes e durante o evento (antecipação), ficou abaixo do normal
depois dele (desmobilização), com efeito concentrado na Região Metropolitana de Belém e no alojamento, e
que o efeito se dissipa em poucos meses.

\section{Contribuição}

A dissertação contribui em três frentes. Analisa um megaevento \textbf{não esportivo}, tipo de evento
raramente estudado (Bronzini, Mocetti e Mongardini, 2019). Estuda uma \textbf{cidade amazônica}, com
infraestrutura hoteleira menor e mais sensível a choques. E usa \textbf{microdados mensais} do emprego
formal em um painel de capitais, estimando o efeito por diferenças em diferenças e controle sintético.

\section{Estrutura da dissertação}

\begin{itemize}
  \item O Capítulo~\ref{cap:referencial} apresenta o referencial teórico.
  \item O Capítulo~\ref{cap:dados} descreve as fontes, a construção da base e o mecanismo.
  \item O Capítulo~\ref{cap:micro}, central, apresenta a estratégia de identificação e os resultados:
        Belém comparada com as outras capitais, por diferenças em diferenças e controle sintético.
  \item O Capítulo~\ref{cap:macro} traz evidência complementar para o Pará, com uma regressão com
        dummies de fase e controles macroeconômicos.
  \item O Capítulo~\ref{cap:conclusao} conclui.
\end{itemize}
